\documentclass[10pt,a4paper]{amsart}
\usepackage[margin=19mm]{geometry}
\usepackage{amsmath,amssymb,mathtools}
\usepackage[scr=rsfs,cal=euler]{mathalfa}
\usepackage{xfrac}
\usepackage[unicode,hidelinks,bookmarks=false]{hyperref}
\newcommand{\ie}{i.e.\ }
\newcommand{\cf}{cf.\ }
\newcommand{\resp}{respectively\ }
\newcommand{\Subsection}{\subsection}
\providecommand{\nobreakdash}{\nobreak\hbox{-}}
\setlength{\emergencystretch}{2em}
\multlinegap0pt
\usepackage{tikz}
\usetikzlibrary{cd}
\usepackage{mathtools}
\usepackage{amssymb}
\newtheorem{corollary}{Corollary}
\newtheorem{lemma}{Lemma}
\newtheorem{proposition}{Proposition}
\def\Kl{\mathrm{Kl}}
\def\PP{\mathbb{P}}
\def\Pa{\mathrm{Pa}}
\def\ZZ{\mathbb{Z}}
\def\bfG{\mathbf{G}}
\def\calO{\mathcal{O}}
\def\calT{\mathcal{T}}
\def\lr{\mathrm{lr}}
\def\rmB{\mathrm{B}}
\def\rmC{\mathrm{C}}
\def\rmH{\mathrm{H}}
\def\rmK{\mathrm{K}}
\def\rmN{\mathrm{N}}
\def\rmR{\mathrm{R}}
\def\rmT{\mathrm{T}}
\def\rmX{\mathrm{X}}
\def\rmZ{\mathrm{Z}}
\def\rmc{\mathrm{c}}
\title{Arithmetic level raising for certain quaternionic unitary Shimura variety}
\author{Haining Wang}
\date{Source: arXiv:2204.06976v4 (CC BY 4.0)}
\begin{document}
\maketitle

\noindent\emph{Source and licence.} Arithmetic level raising for certain quaternionic unitary Shimura variety. Version arXiv:2204.06976v4 (CC BY 4.0), distributed by arXiv under the Creative Commons Attribution 4.0 International licence, http://creativecommons.org/licenses/by/4.0/, as recorded in the arXiv licence metadata for this exact version and shown on the abstract page https://arxiv.org/abs/2204.06976v4. Full licence notice: \texttt{licenses/arxiv-2204.06976-CC-BY-4.0.txt}.\par\smallskip

\noindent\textit{Editorial scope.} Counted unit: the article's proposition lr-matrix (source0.tex lines 2123--2128). The article's complete original proof of it is delivered with it (source0.tex lines 2130--2226, one \texttt{proof} environment of 97 lines). No mathematics is written, repaired or re-derived: the statement and the proof are the article's own lines, in the article's own notation, and no retained line is substituted or re-flowed.  Retained uncounted as the source-local context the proof needs: lines 1032--1053; lines 1724--1732; lines 1774--1787; lines 1831--1853; lines 2015--2024.  Nothing was omitted from any retained span, so the package carries no omission record.  No retained line contains a citation, so the wrapper carries no bibliography and no bibliography entry.  The retained lines contain the article's own diagram code, which the wrapper typesets with the standard tikz packages; no image file is delivered and the package declares no asset dependency.  Editorial formatting only, in addition: an AMS \texttt{amsart} preamble that restates the article's own declarations and macros used by the retained lines, namely the environments \texttt{corollary}, \texttt{lemma}, \texttt{proposition} on separate counters, together with the standard packages those lines need.  The retained environments print in this document as Corollary 1, Lemma 1, Proposition 1 in order of appearance, whereas the article prints the same items with the numbering of its own full document.  The complete native source of the article as distributed in the arXiv e-print for this exact version is bundled unchanged as \texttt{sources/124545/source0.tex}.
\begin{corollary}\label{ss-bullet}
The scheme $\overline{\rmX}^{\bullet\mathrm{ss}}_{\Pa}(\rmB)$ is pure of dimension $2$ whose irreducible components are parametrized by 
the Shimura sets $\rmZ_{\{0\}}(\overline{\rmB})$ and $\rmZ_{\{2\}}(\overline{\rmB})$. 
\begin{enumerate}

\item An irreducible component of $\overline{\rmX}^{\bullet\mathrm{ss}}_{\Pa}(\rmB)$ is isomorphic to $\mathrm{DL}^{\bullet}(\Lambda_{i})$ for some vertex lattice $\Lambda_{i}$ of type $i\in\{0,2\}$ which is isomorphic to the blow-up of $\mathrm{DL}(\Lambda_{i})$ at all the superspecial points contained in $\mathrm{DL}(\Lambda_{i})$.

\item For each pair of vertex lattices $(\Lambda_{0}, \Lambda_{2})$ such that
\begin{equation*}
p\Lambda^{\vee}_{0}\subset p\Lambda^{\vee}_{2}=\Lambda_{2}\subset \Lambda_{0}, 
\end{equation*}
the corresponding irreducible component $\mathrm{DL}^{\bullet}(\Lambda_{0})$ and $\mathrm{DL}^{\bullet}(\Lambda_{2})$ intersects at a projective line $\PP^{1}$.

\item For each vertex lattice $\Lambda_{1}$ of type $1$ such that $\Lambda_{1}\subset \Lambda_{0}$ for a vertex lattice $\Lambda_{0}$ of type $0$, $\mathrm{DL}^{\bullet}(\Lambda_{0})$ and  $\rmC^{\bullet}(\Lambda_{1})$ intersects at a projective line $\PP^{1}$.

\item For each vertex lattice $\Lambda_{1}$ of type $1$ such that $\Lambda_{2}\subset \Lambda_{1}$ for a vertex lattice $\Lambda_{2}$ of type $2$, $\mathrm{DL}^{\bullet}(\Lambda_{2})$ and $\rmC^{\bullet}(\Lambda_{1})$ intersects at a projective line $\PP^{1}$.

\item For each pair $(\Lambda_{0}, \Lambda^{\prime}_{0})$ of vertex lattices of type $0$ such that $\Lambda_{1}=\Lambda_{0}\cap\Lambda^{\prime}_{0}$ is a vertex lattice of type $1$, the irreducible components $\mathrm{DL}^{\bullet}(\Lambda_{0})$ and $\mathrm{DL}^{\bullet}(\Lambda^{\prime}_{0})$ has trivial intersection unless $\Lambda_{0}=\Lambda^{\prime}_{0}$.

\item For each pair $(\Lambda_{2}, \Lambda^{\prime}_{2})$ of vertex lattices of type $2$ such that $\Lambda_{1}=\Lambda_{2}+\Lambda^{\prime}_{2}$ is a vertex lattice of type $1$, the irreducible components $\mathrm{DL}^{\bullet}(\Lambda_{2})$ and $\mathrm{DL}^{\bullet}(\Lambda^{\prime}_{2})$ has trivial intersection unless $\Lambda_{2}=\Lambda^{\prime}_{2}$.
\end{enumerate}
\end{corollary}

\begin{corollary}\label{intersection}
We have the following isomorphism
\begin{equation*}
\begin{aligned}
\mathrm{IH}^{n}_{(\rmc)}(\rmX_{\overline{s}}, \Lambda)&\cong\mathrm{Im}(\mathrm{H}^{n}_{(\rmc)}(\rmX_{\overline{s}}, \Lambda)\rightarrow \mathrm{H}^{n}_{(\rmc)}(\rmX_{\overline{s}}, j_{\Sigma\ast}\Lambda))\\
&\cong\mathrm{H}^{n}_{(\rmc)}(\rmX_{\overline{s}}, \Lambda)/\bigoplus\limits_{\sigma\in\Sigma}\rmH^{n}_{\{\sigma\}}(\rmX_{\overline{s}}, \Lambda).\\
\end{aligned}
\end{equation*}
\end{corollary}

\begin{lemma}\label{quadric-class}
The cohomology group of the quadric $\rmC_{\sigma}$ is given by the formulas below.
\begin{enumerate}
\item If $n-1=2r-1$ is odd, then we have 
\begin{equation*}
\rmH^{i}(\rmC_{\sigma}, \Lambda(i/2))=\begin{cases}\Lambda\eta^{i/2} & \text{if $i$ even, $0\leq i\leq 2(n-1)$;}\\ 0 & \text{otherwise.}\\ \end{cases}
\end{equation*}
\item If $n-1=2r-2$ is even, then we have 
\begin{equation*}
\rmH^{i}(\rmC_{\sigma}, \Lambda(i/2))=\begin{cases}\Lambda\eta^{i/2} & \text{if $i$ even, $0\leq i\leq 2(n-1), i\neq n-1$;}\\   \Lambda\theta_{1}\oplus\Lambda\theta_{2} & \text{if $i=n-1$;}\\0 & \text{otherwise.}\\ \end{cases}
\end{equation*}
\end{enumerate}
Here $\eta\in \rmH^{2}(\rmC_{\sigma},\Lambda(1))$ is the class of the hyperplane section of $\rmC_{\sigma}$ and $\theta_{1}$ and $\theta_{2}$ are the cycle classes of the two generatrices of $\rmC_{\sigma}$ such that $\theta_{1}+\theta_{2}=\eta^{(n-1)/2}$ holds.
\end{lemma}

\begin{proposition}\label{alpha-beta}
We have the following descriptions of the maps $\alpha_{\rmc}$ and $\beta$.
\begin{enumerate}
\item The space of vanishing cycles  $\bigoplus\limits_{\sigma\in\Sigma}\rmR^{n}\Phi_{\sigma}(\Lambda)(r)$ and the map $\alpha_{\rmc}$ inserts in the following commutative diagram 
\begin{equation*}
\begin{tikzcd}
\bigoplus\limits_{\sigma\in\Sigma}\rmH^{n+1}(\rmC_{\sigma}, \Lambda(r)) \arrow[r, "\sim"]                 &       \bigoplus\limits_{\sigma\in\Sigma}\rmH^{n+1}(\rmC_{\sigma}, \Lambda(r))  \\
\bigoplus\limits_{\sigma\in\Sigma}\rmH^{n-1}(\rmC_{\sigma}, \Lambda(r-1)) \arrow[r] \arrow[u, "\cup \eta"'] & \rmH^{n+1}_{\rmc}({\rmX}^{\bullet}_{\overline{s}}, \Lambda(r)) \arrow[u] \\
\bigoplus\limits_{\sigma\in\Sigma}\rmR^{n}\Phi_{\sigma}(\Lambda)(r) \arrow[r, "\alpha_{\rmc}(r)"] \arrow[u]       & \rmH^{n+1}_{\rmc}(\rmX^{\square}_{\overline{s}}, \Lambda(r)) \arrow[u]
\end{tikzcd}
\end{equation*}
where the left and right columns are exact.
\item The space of co-vanishing cycles   $\bigoplus\limits_{\sigma\in\Sigma}\rmH^{n}_{\{\sigma\}}(\rmX_{\overline{s}}, \rmR\Psi(\Lambda)(r-1))$ and the map $\beta$ inserts in the following commutative diagram 
\begin{equation*}
\begin{tikzcd}
\bigoplus\limits_{\sigma\in\Sigma}\rmH^{n}_{\{\sigma\}}(\rmX_{\overline{s}}, \rmR\Psi(\Lambda)(r-1))              &\mathrm{H}^{n-1}(\rmX^{\square}_{\overline{s}}, \Lambda(r-1))\arrow[l, "\beta(r-1)"]           \\
\bigoplus\limits_{\sigma\in\Sigma}\rmH^{n-1}(\rmC_{\sigma}, \Lambda(r-1)) \arrow[u] &   \mathrm{H}^{n-1}({\rmX}^{\bullet}_{\overline{s}}, \Lambda(r-1)) \arrow[l] \arrow[u] \\
\bigoplus\limits_{\sigma\in\Sigma}\rmH^{n-3}(\rmC_{\sigma}, \Lambda(r-2)) \arrow[u, "\cup\eta"]       &  \bigoplus\limits_{\sigma\in\Sigma}\rmH^{n-3}(\rmC_{\sigma}, \Lambda(r-2))\arrow[l,"\sim"] \arrow[u]
\end{tikzcd}
\end{equation*}
where the left and right columns are exact.
\end{enumerate}
\end{proposition}

\begin{equation}\label{level-raise-diagram}
\begin{tikzcd}
 \calO_{\lambda}[\rmZ_{\{0\}}(\overline{\rmB})] \arrow[rd, "\mathrm{inc}^{\{0\}}_{!}"'] &                                    &  \calO_{\lambda}[\rmZ_{\{2\}}(\overline{\rmB})]\arrow[ld, "\mathrm{inc}^{\{2\}}_{!}"] \\
                   & \rmH^{2}(\overline{\rmX}^{\square}_{\Pa}(\rmB), \calO_{\lambda}(1)) \arrow[d, "\beta(1)"]                   &                   \\
                   & \bigoplus\limits_{\sigma\in \rmZ_{\{1\}}(\overline{\rmB})}\rmH^{3}_{\{\sigma\}}(\overline{\rmX}_{\Pa}(\rmB), \rmR\Psi(\calO_{\lambda})(1)) \arrow[d, "\mathrm{N}^{-1}_{\Sigma}"]                   &                   \\
                   & \bigoplus\limits_{\sigma\in \rmZ_{\{1\}}(\overline{\rmB})}\rmR^{3}\Phi_{\sigma}(\calO_{\lambda})(2)) \arrow[d, "\alpha_{\rmc}(2)"]                   &                   \\
                   &  \rmH^{4}_{\rmc}(\overline{\rmX}^{\square}_{\Pa}(\rmB), \calO_{\lambda}(2)) \arrow[ld, "\mathrm{inc}^{\ast}_{\{0\}}"'] \arrow[rd, "\mathrm{inc}^{\ast}_{\{2\}}"] &                   \\
 \calO_{\lambda}[\rmZ_{\{0\}}(\overline{\rmB})]            &                                    &   \calO_{\lambda}[\rmZ_{\{2\}}(\overline{\rmB})]            
\end{tikzcd}
\end{equation}

\begin{proposition}\label{lr-matrix}
The level raising matrix $\calT_{\mathrm{lr}}$ is given by
\begin{equation*}
\calT_{\mathrm{lr}}=\begin{pmatrix} &-(\rmT^{-1}_{p,0}\rmT_{p,1}+(p+1)(p^{2}+1)) & (p+1)\rmT_{02}\\ &(p+1)\rmT_{20}& -(\rmT^{-1}_{p,0}\rmT_{p,1}+(p+1)(p^{2}+1))\\ \end{pmatrix}.
\end{equation*}
\end{proposition}

\begin{proof}
To begin the proof, note that we can identify both the space of vanishing cycles
\begin{equation*}
\bigoplus\limits_{\sigma\in \rmZ_{\{1\}}(\overline{\rmB})}\rmH^{3}_{\{\sigma\}}(\overline{\rmX}_{\Pa}(\rmB), \rmR\Psi(\calO_{\lambda}(1))) 
\end{equation*}
and the space of co-vanishing cycles
\begin{equation*}
\bigoplus\limits_{\sigma\in \rmZ_{\{1\}}(\overline{\rmB})}\rmR^{3}\Phi_{\sigma}(\calO_{\lambda}(2)) 
\end{equation*}
with $\calO_{\lambda}[\rmZ_{\Pa}(\overline{\rmB})]$ ignoring the Galois actions using the generators $\eta^{\sigma}_{\bullet}\in\rmH^{3}_{\{\sigma\}}(\overline{\rmX}_{\Pa}(\rmB), \rmR\Psi(\calO_{\lambda}(1)))$ and $\eta^{\bullet}_{\sigma}\in \rmR^{3}\Phi_{\sigma}(\calO_{\lambda}(2))$ in Lemma \ref{intersection}. The level raising diagram \eqref{level-raise-diagram} reduces to 
\begin{equation}
\begin{tikzcd}
 \calO_{\lambda}[\rmZ_{\{0\}}(\overline{\rmB})] \arrow[rd, "\mathrm{inc}^{\{0\}}_{!}"'] &                                    &  \calO_{\lambda}[\rmZ_{\{2\}}(\overline{\rmB})]\arrow[ld, "\mathrm{inc}^{\{2\}}_{!}"] \\
                   & \rmH^{2}(\overline{\rmX}^{\square}_{\Pa}(\rmB), \calO_{\lambda}(1)) \arrow[d, "\beta(1)"]                   &                   \\
                   & \calO_{\lambda}[\rmZ_{\Pa}(\overline{\rmB})] \arrow[d, "\rmN^{-1}_{\Sigma}"]   &\\
                   & \calO_{\lambda}[\rmZ_{\Pa}(\overline{\rmB})] \arrow[d, "\alpha_{\rmc}(2)"]                   &                   \\
                   &  \rmH^{4}_{\rmc}(\overline{\rmX}^{\square}_{\Pa}(\rmB), \calO_{\lambda}(2)) \arrow[ld, "\mathrm{inc}^{\ast}_{\{0\}}"'] \arrow[rd, "\mathrm{inc}^{\ast}_{\{2\}}"] &                   \\
 \calO_{\lambda}[\rmZ_{\{0\}}(\overline{\rmB})]            &                                    &   \calO_{\lambda}[\rmZ_{\{2\}}(\overline{\rmB})].            
\end{tikzcd}
\end{equation}
This justifies the appearance of the Hecke operators appearing in the level raising matrix.
By Proposition \ref{alpha-beta}, the map $\alpha_{\rmc}(2)$ can be realized as a Gysin map and $\beta(1)$ can be realized as a restriction map. In fact, we have the following  commutative diagram where all the rows are exact:
\begin{equation*}
\begin{tikzcd}
\rmH^{2}(\overline{\rmX}^{\square}_{\Pa}(\rmB), \calO_{\lambda}(1)) \arrow[d] & \rmH^{2}(\overline{\rmX}^{\bullet}_{\Pa}(\rmB), \calO_{\lambda}(1)) \arrow[l] \arrow[d] & \bigoplus\limits_{\sigma\in \rmZ_{\Pa}(\overline{\rmB})}\rmH^{0}(\rmC_{\sigma}, \calO_{\lambda}) \arrow[l] \arrow[d, equal] \\
\bigoplus\limits_{\sigma\in \rmZ_{\Pa}(\overline{\rmB})}\rmH^{3}_{\{\sigma\}}(\overline{\rmX}_{\Pa}(\rmB), \calO_{\lambda}(1))  \arrow[d, "\rmN^{-1}_{\Sigma}"]           &  \bigoplus\limits_{\sigma\in \rmZ_{\Pa}(\overline{\rmB})}\rmH^{2}(\rmC_{\sigma}, \calO_{\lambda}(1))  \arrow[l]           &  \bigoplus\limits_{\sigma\in \rmZ_{\Pa}(\overline{\rmB})}\rmH^{0}(\rmC_{\sigma}, \calO_{\lambda}) \arrow[l]                                \\
\bigoplus\limits_{\sigma\in \rmZ_{\Pa}(\overline{\rmB})}\rmR^{3}\Phi_{\sigma}(\calO_{\lambda}(2))  \arrow[r] \arrow[d] & \bigoplus\limits_{\sigma\in \rmZ_{\Pa}(\overline{\rmB})}\rmH^{2}(\rmC_{\sigma}, \calO_{\lambda}(1))  \arrow[r] \arrow[d] & \bigoplus\limits_{\sigma\in \rmZ_{\Pa}(\overline{\rmB})}\rmH^{4}(\rmC_{\sigma}, \calO_{\lambda}(2))  \arrow[d, equal]           \\
\rmH^{4}_{\rmc}(\overline{\rmX}^{\square}_{\Pa}(\rmB), \calO_{\lambda}(2)) \arrow[r]           & \rmH^{4}_{\rmc}(\overline{\rmX}^{\bullet}_{\Pa}(\rmB), \calO_{\lambda}(2))\arrow[r]           &\bigoplus\limits_{\sigma\in \rmZ_{\Pa}(\overline{\rmB})}\rmH^{4}(\rmC_{\sigma}, \calO_{\lambda}(2)).                                   
\end{tikzcd}
\end{equation*}

Let $\overline{\rmX}^{\bullet\mathrm{ss}}_{\Pa}(\rmB)$ be the supersingular locus of $\overline{\rmX}^{\bullet}_{\Pa}(\rmB)$. We know that it has three types of irreducible components parametrized by $\rmZ_{\{i\}}(\overline{\rmB})$ for $i\in\{0,2\}$ and $\rmZ_{\Pa}(\overline{\rmB})$. We denote by $\overline{\rmX}^{\bullet\mathrm{ss}.n}_{\Pa, \{i\}}(\rmB)$ the disjoint union of the irreducible components of $\overline{\rmX}^{\bullet\mathrm{ss}}_{\Pa}(\rmB)$ parametrized by $\rmZ_{\{i\}}(\overline{\rmB})$ for $i\in\{0,2\}$. 

Therefore, for $i\in\{0,2\}$, we have the natural Gysin morphism
\begin{equation}\label{gys-sm}
\mathrm{inc}^{\{i\}}_{!}: \calO_{\lambda}[\rmZ_{\{i\}}(\overline{\rmB})]\cong\rmH^{0}(\overline{\rmX}^{\bullet\mathrm{ss}.n}_{\Pa, \{i\}}(\rmB), \calO_{\lambda})\rightarrow  \rmH^{2}(\overline{\rmX}^{\bullet}_{\Pa}(\rmB), \calO_{\lambda}(1))
\end{equation}
and the natural restriction morphism
\begin{equation}\label{res-sm}
\mathrm{inc}^{\ast}_{\{i\}}: \rmH^{4}_{\rmc}(\overline{\rmX}^{\bullet}_{\Pa}(\rmB), \calO_{\lambda}(2))\rightarrow  \rmH^{4}_{\rmc}(\overline{\rmX}^{\bullet\mathrm{ss}.n}_{\Pa, \{i\}}(\rmB), \calO_{\lambda}(2)) \cong \calO_{\lambda}[\rmZ_{\{i\}}(\overline{\rmB})].
\end{equation}
By the commutativity of the above diagram, to compute the entry $\calT_{\mathrm{lr},\{ij\}}$ for $i, j\in\{0, 2\}$, it suffices to understand the compotite map given by
\begin{equation*}
\begin{aligned}
&\calO_{\lambda}[\rmZ_{\{i\}}(\overline{\rmB})]\xrightarrow{\mathrm{inc}^{\{i\}}_{!}}  \rmH^{2}({\overline{\rmX}}^{\bullet}_{\Pa}(\rmB), \calO_{\lambda}(1))\rightarrow \bigoplus\limits_{\sigma\in \rmZ_{\Pa}(\overline{\rmB})}\rmH^{2}(\rmC_{\sigma}, \calO_{\lambda}(1))\rightarrow \bigoplus\limits_{\sigma\in \rmZ_{\Pa}(\overline{\rmB})}\rmH^{3}_{\{\sigma\}}(\overline{\rmX}_{\Pa}(\rmB), \calO_{\lambda}(1))\\
&\xrightarrow{\rmN^{-1}_{\Sigma}}\bigoplus\limits_{\sigma\in \rmZ_{\Pa}(\overline{\rmB})}\rmR^{3}\Phi_{\sigma}(\calO_{\lambda}(2)) \rightarrow \bigoplus\limits_{\sigma\in \rmZ_{\Pa}(\overline{\rmB})}\rmH^{2}(\rmC_{\sigma}, \calO_{\lambda}(1))\rightarrow \rmH^{4}_{\rmc}({\overline{\rmX}}^{\bullet}_{\Pa}(\rmB), \calO_{\lambda}(2))\xrightarrow{\mathrm{inc}^{\ast}_{\{j\}}} \calO_{\lambda}[\rmZ_{\{j\}}(\overline{\rmB})].
\end{aligned}
\end{equation*}
Using Lemma \ref{quadric-class}, \ref{intersection} and Corollary \ref{ss-bullet}, we arrive at the following term by term computation.

\subsubsection{$\calT_{\mathrm{lr},\{00\}}$} Note that $\rmZ_{\{i\}}(\overline{\rmB})$ corresponds to the set of vertex lattices of type $i$ for $i\in\{0,  2\}$ while $\rmZ_{\Pa}(\overline{\rmB})$ corresponds to the set of vertex lattices of type $1$.  Therefore to understand the map $\calT_{\lr, \{00\}}=\mathrm{inc}^{\ast}_{\{0\}} \circ\alpha_{\rmc}(2)\circ\rmN^{-1}_{\Sigma}\circ\beta(1)\circ\mathrm{inc}^{\{0\}}_{!}$, we need to understand the possible relative position of vertex lattices $\Lambda_{0}$, $\Lambda_{\Pa}$ and $\Lambda^{\prime}_{0}$ of type $0$, $1$ and $0$. It not hard to see that the only possible situation is that $\Lambda_{\Pa}\subset^{1} \Lambda_{0}$ and  $\Lambda_{\Pa}\subset^{1} \Lambda^{\prime}_{0}$. There are two possible cases.

Case $(1)$: $\Lambda_{0}=\Lambda^{\prime}_{0}$. Then we have $p\Lambda^{\vee}_{0}\subset^{1}p\Lambda^{\vee}_{\Pa}\subset^{2} \Lambda_{\Pa}\subset ^{1}\Lambda_{0}$. Giving such a $\Lambda_{\Pa}$ is equivalent to giving a complete flag in $\Lambda_{0}/p\Lambda_{0}$. The number of such flags is given by 
\begin{equation*}
[\bfG(\ZZ_{p}): \Kl]=(p+1)(p^{2}+1) 
\end{equation*}
where $\Kl$ is the Klingen parahoric.
 
Case $(2)$: $\Lambda_{0}\neq\Lambda^{\prime}_{0}$. Then we have $\Lambda_{0}\cap\Lambda^{\prime}_{0}=\Lambda_{\Pa}$ and the relative position between $\Lambda_{0}$ and $\Lambda^{\prime}_{0}$ is given by
\begin{equation*}
\rmK_{\{0\}}\begin{pmatrix}p&&&\\&1&&\\&&1&\\&&&p^{-1}\\\end{pmatrix}\rmK_{\{0\}}.
\end{equation*}
This operator is the same as $\rmT^{-1}_{p,0}\rmT_{p,1}$ by identifying $\rmK_{\{0\}}$ with the hyperspecial group $\bfG(\ZZ_{p})$. Therefore the entry $\calT_{\mathrm{lr}, \{00\}}$ is given by  
\begin{equation*}
\calT_{\mathrm{lr}, \{00\}}=-(\rmT^{-1}_{p,0}\rmT_{p,1}+(p+1)(p^{2}+1)). 
\end{equation*}
Here the sign $-1$ comes from the formula in Lemma \ref{intersection} $(3)$.

\subsubsection{$\calT_{\mathrm{lr}, \{22\}}$} To understand the map $\calT_{\lr, \{22\}}=\mathrm{inc}^{\ast}_{\{2\}} \circ\alpha_{\rmc}(2)\circ\rmN^{-1}_{\Sigma}\circ\beta(1)\circ\mathrm{inc}^{\{2\}}_{!}$, we need to understand the possible relative position of vertex lattices $\Lambda_{2}$, $\Lambda_{\Pa}$ and $\Lambda^{\prime}_{2}$ of type $2$, $1$ and $2$ respectively. It not hard to see that the only possible way is that $\Lambda_{2}\subset^{1}\Lambda_{\Pa}$ and  $\Lambda^{\prime}_{2}\subset^{1}\Lambda_{\Pa}$. The rest of the analysis is completely the same as in the case of $\calT_{\lr,\{00\}}$. Once we identify $\rmK_{\{2\}}$ with the hyperspecial group $\bfG(\ZZ_{p})$, the entry $\calT_{\lr, \{22\}}$ is given by 
\begin{equation*}
\calT_{\lr, \{22\}}=-(\rmT^{-1}_{p,0}\rmT_{p,1}+(p+1)(p^{2}+1)). 
\end{equation*}
Here the sign $-1$ comes from the formula in Lemma \ref{intersection} $(3)$.

\subsubsection{$\calT_{\mathrm{lr}, \{02\}}$} To understand the map $\calT_{\lr, \{02\}}=\mathrm{inc}^{\ast}_{\{2\}} \circ\alpha_{\rmc}(2)\circ\rmN^{-1}_{\Sigma}\circ\beta(1)\circ\mathrm{inc}^{\{0\}}_{!}$, we need to understand the possible relative position of vertex lattices $\Lambda_{0}$, $\Lambda_{\Pa}$ and $\Lambda^{\prime}_{2}$ of type $0$, $1$ and $2$. We find that the only possible situation is that
\begin{equation*}
\Lambda_{2}\subset^{1}\Lambda_{\Pa}\subset^{1}\Lambda_{0}=\Lambda^{\vee}_{0}\subset^{1}\Lambda^{\vee}_{\Pa}\subset^{1}\Lambda^{\vee}_{2}.
\end{equation*}
It follows each such $\Lambda_{\Pa}$ corresponds to a line in the two dimensional space $\Lambda_{0}/\Lambda_{2}$. 

Therefore the entry $\calT_{\lr, \{02\}}$ is given by
\begin{equation*}
\calT_{\lr, \{02\}}=(p+1)\mathrm{char}(\rmK_{\{0\}}\begin{pmatrix}p&&&\\&p&&\\&&1&\\&&&1 \end{pmatrix}\rmK_{\{2\}})=(p+1)\rmT_{02}.\
\end{equation*}

\subsubsection{$\calT_{\mathrm{lr}, \{20\}}$} To understand the map $\calT_{\mathrm{lr}, \{20\}}=\mathrm{inc}^{\ast}_{\{0\}} \circ\alpha_{\rmc}(2)\circ\rmN^{-1}_{\Sigma}\circ\beta(1)\circ\mathrm{inc}^{\{2\}}_{!}$, we need to understand the possible relative positions of the vertex lattices $\Lambda_{0}$, $\Lambda_{\Pa}$ and $\Lambda^{\prime}_{2}$ of type $0$, $1$ and $2$. We find that the only possible way is that
\begin{equation*}
\Lambda_{2}\subset^{1}\Lambda_{\Pa}\subset^{1}\Lambda_{0}=\Lambda^{\vee}_{0}\subset^{1}\Lambda^{\vee}_{\Pa}\subset^{1}\Lambda^{\vee}_{2}.
\end{equation*}
Each $\Lambda_{\Pa}$ corresponds to a line in the two dimensional space $\Lambda_{0}/\Lambda_{2}$. 

Therefore the entry $\calT_{\mathrm{lr},\{20\}}$ is given by
\begin{equation*}
\calT_{\mathrm{lr},\{20\}}=(p+1)\mathrm{char}(\rmK_{\{0\}}\begin{pmatrix}p^{-1}&&&\\&p^{-1}&&\\&&1&\\&&&1 \end{pmatrix}\rmK_{\{2\}})=(p+1)\rmT_{02}.\
\end{equation*}
This finishes the proof of Proposition \ref{lr-matrix}.
\end{proof}

\end{document}
